\chapter{问题定义与总体框架设计}

多轮对话检索基准的构建涉及文档语料库的组织、对话实例的结构化表示、标注集的质量保障以及评估指标的合理设计等多个紧密关联的问题。本章首先对多轮对话检索任务的核心概念进行严格的形式化定义（第\ref{sec:formal_definition}节），随后从理论层面系统分析现有基准数据集中标注噪声与标注稀疏性两类核心缺陷的形式化描述及其对评估结果的影响机制（第\ref{sec:defect_analysis}节），最后在此基础上提出CRB-Suite框架的总体架构设计，阐明各模块之间的逻辑关系与协作机制（第\ref{sec:framework_overview}节）。

\section{多轮对话检索基准的形式化定义}
\label{sec:formal_definition}

\subsection{文档语料库与对话实例}

多轮对话检索任务的核心要素包括一个文档语料库和一组结构化的对话评估实例。我们对这些要素进行如下形式化定义。

文档语料库。设文档语料库为 $\mathcal{D} = \{d_1, d_2, \dots, d_N\}$，其中 $N = |\mathcal{D}|$ 表示语料库中文档的总数，每个 $d_j$（$j = 1, 2, \dots, N$）表示一个独立的文档单元（或段落）。在实际的检索增强生成（RAG）系统中，原始文档通常需要经过分块（chunking）处理，将长文档切分为适合检索模型处理的较短段落。因此，此处的"文档"$d_j$ 既可以指一篇完整的文章，也可以指经过分块后的一个文本段落——其具体粒度取决于实际应用场景和评估需求。文档语料库的规模 $N$ 直接决定了检索任务的搜索空间大小，从而影响检索难度。现有基准数据集的语料库规模从数百篇（如Doc2Dial~\citep{feng-etal-2020-doc2dial}的488篇文档）到数千万篇（如TopiOCQA~\citep{adlakha-etal-2022-topiocqa}的约2,000万篇Wikipedia文档）不等。本文则选用 FineWeb-Edu~\citep{penedo2024finewebedufineweb} 作为底层语料基座，该数据集源自 Common Crawl 2024 网页快照并经 LLM 教育价值分类器过滤，约 \TBD{FineWeb-Edu 规模} 的规模与高教育值内容为大规模对话合成提供了时效性与覆盖度兼备的知识源。

多轮对话实例。一段多轮对话 $\mathcal{C}$ 由 $T$ 个连续的交互轮次组成。其简单公式表示如下所示：
\begin{equation}
    \mathcal{C} = \{(q_1, a_1), (q_2, a_2), \dots, (q_T, a_T)\}
\end{equation}
其中 $q_i$ 表示第 $i$ 轮用户提出的查询，$a_i$ 表示系统（或标注者）提供的回答。多轮对话的核心特征在于，每一轮的查询 $q_i$ 可能依赖于之前的对话上下文，包含指代消解（如"他是谁？"中的"他"指代前文提及的某个实体）、省略补全（如"那第二个呢？"省略了前文讨论的完整主题）、以及话题切换等复杂的语言现象。

对话历史。对于第 $i$ 轮查询 $q_i$，其对应的对话历史定义为之前所有轮次的查询-回答对的有序序列表示如下：
\begin{equation}
    H_i = \{(q_1, a_1), (q_2, a_2), \dots, (q_{i-1}, a_{i-1})\}
\end{equation}
当 $i = 1$ 时，$H_1 = \emptyset$，即第一轮查询没有历史上下文。对话历史 $H_i$ 提供了理解当前查询 $q_i$ 所必需的上下文信息。在多轮对话检索中，检索模型需要综合考虑当前查询 $q_i$ 和对话历史 $H_i$，以准确理解用户的信息需求。

评估实例。每一个评估实例对应对话中的某一个一个特定轮次 $i$，表示为一个三元组：
\begin{equation}
    \mathcal{E}_i = (H_i, q_i, G_i)
\end{equation}
其中 $G_i \subseteq \mathcal{D}$ 是标注者为回答查询 $q_i$ 而标注的\textbf{金标文档集}（ground-truth document set），即被认为包含回答 $q_i$ 所需信息的文档集合。检索系统的任务是：给定对话历史 $H_i$ 和当前查询 $q_i$，从语料库 $\mathcal{D}$ 中检索出一个排序的候选文档列表 $\mathcal{R}_i = [r_1, r_2, \dots, r_K]$，其中 $K$ 为检索预算（即返回文档数），使得排序靠前的文档尽可能多地覆盖金标文档集 $G_i$ 中的文档。

\subsection{检索性能评估指标}

为了量化检索系统的性能，研究者通常采用一系列标准化的信息检索评估指标。根据是否考虑检索结果的排序位置，这些指标可分为顺序无关指标和顺序敏感指标两大类。

召回率（Recall@$K$）。召回率是最直观的顺序无关指标，衡量检索系统在返回 $K$ 个文档的预算内，成功覆盖了多少比例的金标文档。其定义为：
\begin{equation}
    \text{Recall@}K = \frac{|G_i \cap \mathcal{R}_i^K|}{|G_i|}
\end{equation}
其中 $\mathcal{R}_i^K$ 表示检索结果列表 $\mathcal{R}_i$ 中排名前 $K$ 的文档集合。Recall@$K$ 仅关注金标文档是否被检索到，而不考虑其在结果列表中的具体排名位置。当 $|G_i| = 1$（即每个查询仅有一个金标文档）时，Recall@$K$ 退化为一个二值指标——命中为1，未命中为0。该指标的优势在于高度可解释性，但缺点是无法区分金标文档排在第1位还是第$K$位的情况。

平均倒数排名（MRR@$K$）。MRR是一种顺序敏感指标~\citep{voorhees1999trec}，关注第一个相关文档出现的位置。对于单个查询 $q_i$，倒数排名定义为：
\begin{equation}
    \text{RR}_i = \frac{1}{\text{rank}(d^*_i)}
\end{equation}
其中 $\text{rank}(d^*_i)$ 表示第一个金标文档在检索结果列表 $\mathcal{R}_i$ 中的排名位置。如果在前 $K$ 个结果中没有出现任何金标文档，则 $\text{RR}_i = 0$。MRR@$K$ 是所有查询的倒数排名的平均值：
\begin{equation}
    \text{MRR@}K = \frac{1}{|\mathcal{Q}|} \sum_{i=1}^{|\mathcal{Q}|} \text{RR}_i
\end{equation}
其中 $\mathcal{Q}$ 为所有评估查询的集合。MRR适用于用户只关心第一个正确结果的场景，能够有效区分将金标文档排在首位与排在末位的系统。

归一化折损累计增益（NDCG@$K$）。NDCG（Normalized Discounted Cumulative Gain）由Järvelin和Kekäläinen~\citep{jarvelin2002cumulated}提出，是信息检索领域最广泛使用的顺序敏感指标之一。它不仅考虑相关文档是否被检索到，还考虑其在结果列表中的排名位置，对排名越靠前的相关文档给予越高的权重。具体地，折损累计增益（DCG）定义为：
\begin{equation}
    \text{DCG@}K = \sum_{j=1}^{K} \frac{\text{rel}(r_j)}{\log_2(j + 1)}
\end{equation}
其中 $\text{rel}(r_j)$ 表示排名第 $j$ 位的文档 $r_j$ 的相关性分数。在二值相关性设定下，$\text{rel}(r_j) = 1$ 若 $r_j \in G_i$，否则 $\text{rel}(r_j) = 0$。分母中的对数项起到"折损"作用——位置越靠后，同一相关文档的贡献越小，这符合用户浏览行为的直觉：用户更倾向于关注排名靠前的结果。理想折损累计增益（IDCG）是在最优排序下的DCG值，即将所有金标文档排在最前面时的得分。NDCG通过DCG与IDCG的比值实现归一化：
\begin{equation}
    \text{NDCG@}K = \frac{\text{DCG@}K}{\text{IDCG@}K}
\end{equation}

NDCG@$K$ 的取值范围为 $[0, 1]$，值越大表示检索结果的排序质量越高。

指标间的互补关系。上述三种指标从不同角度刻画检索系统的性能。Recall@$K$ 衡量"是否找到了"，MRR@$K$ 衡量"第一个相关结果排得多靠前"，NDCG@$K$ 则综合衡量"所有相关结果的排序质量"。在实际评估中，通常同时报告多种指标以获得全面的性能画像。本文在后续实验中主要以Recall@$K$（$K=5, 20$）作为核心指标以保证可解释性，同时提供NDCG@20和MRR@20作为补充参考。

\subsection{对话检索的关键技术挑战}

从上述形式化定义出发，多轮对话检索相较于传统的单轮检索面临以下关键技术挑战：

上下文依赖建模。在多轮对话中，当前查询 $q_i$ 的语义理解高度依赖于对话历史 $H_i$。检索模型需要将 $H_i$ 和 $q_i$ 联合编码为一个有效的查询表示，以准确捕捉用户的信息需求。现有方法主要分为两大类~\citep{mo2025survey}：（1）\textbf{查询改写}，利用大语言模型将依赖上下文的 $q_i$ 改写为一个独立的、去上下文化的查询 $q'_i$，然后使用标准的单轮检索模型进行检索；（2）\textbf{对话密集检索}，训练专门的编码器来直接建模包含历史上下文的查询表示。两种方法各有优劣，但均需要高质量的基准数据集来进行训练和评估。

话题切换处理。真实的人类信息寻求行为并非线性的。Spink等人的研究~\citep{spink2002multitasking}表明，用户在网络搜索过程中平均每2.11轮就会切换话题。话题切换可以是平滑的语义过渡（软切换），也可以是完全无关的突然跳转（硬切换）。当话题发生切换时，对话历史 $H_i$ 中来自之前话题的信息可能对当前查询 $q_i$ 产生干扰，形成噪声上下文。检索模型必须具备区分当前话题与历史噪声的能力，这对模型的上下文理解能力提出了极高的要求。

长上下文处理。随着对话轮次 $T$ 的增加，对话历史 $H_i$ 的长度持续增长。在生产环境中，由于大语言模型经过RLHF训练后倾向于生成冗长的回复~\citep{ouyang2022traininglanguagemodelsfollow}，每一轮的回答 $a_i$ 可能包含数百个token，导致对话历史迅速膨胀。这不仅增加了检索模型的计算负担，更带来了信息过载的挑战——检索模型需要从大量的历史信息中提取与当前查询真正相关的线索，同时忽略无关的冗余内容。


\section{现有基准的核心缺陷分析}
\label{sec:defect_analysis}

在上一节的形式化框架下，我们可以更加精确地描述现有基准数据集中存在的两类核心缺陷。这两类缺陷从根本上影响了评估结果的可靠性，使得在这些基准上获得的性能指标可能无法真实反映检索系统的实际能力。

\subsection{标注噪声}

定义。对于一个评估实例 $\mathcal{E}_i = (H_i, q_i, G_i)$，设 $\hat{G}_i \subseteq \mathcal{D}$ 为在语料库 $\mathcal{D}$ 中\textbf{真正能够}回答查询 $q_i$ 的所有文档的集合（即"理想金标集"）。\textbf{标注噪声}（Annotation Noise）发生在标注的金标文档集中包含了实际上无法回答该查询的文档的情况，即：
\begin{equation}
    \mathcal{N}_i = G_i \setminus \hat{G}_i \neq \emptyset
\end{equation}
其中 $\mathcal{N}_i$ 为第 $i$ 个实例的噪声文档集。换言之，存在文档 $d \in G_i$ 使得 $d \notin \hat{G}_i$，即该文档被错误地标注为"相关"，但实际上并不包含回答 $q_i$ 所需的信息。

成因分析。标注噪声的产生主要有以下几种原因：（1）\textbf{标注者理解偏差}，标注者对查询意图的理解与实际语义存在偏差，导致将主题相关但未直接回答问题的文档标注为金标；（2）\textbf{自动化规则的粗粒度匹配}，在自动化合成方法中，基于关键词匹配或文档结构的规则可能将表面相关但语义不匹配的文档纳入金标集。例如，CORAL~\citep{coral2025}使用正文中的引用链接来确定金标文档，但被引用的文档可能与当前查询的具体问题无关；（3）\textbf{跨轮次标注错位}，在多轮对话中，标注者可能将前一轮查询的金标文档错误地延续到当前轮次，即使当前查询已经切换了话题。

对评估结果的影响。标注噪声导致\textbf{模型性能的过高估计}。当金标集 $G_i$ 中包含噪声文档 $d \in \mathcal{N}_i$ 时，一个"好"的检索器可能正确地将该文档排在较低位置（因为它确实不太相关），但评估指标却因此给予惩罚。反之，一个"差"的检索器如果恰好检索到了这些噪声文档，反而获得了不应有的奖励。这扭曲了评估指标的诊断能力，使得指标分数无法准确反映检索质量的真实差异。

\subsection{标注稀疏性}

定义。\textbf{标注稀疏性}（Annotation Sparsity）是指在语料库中存在能够回答查询 $q_i$ 的有效文档，但这些文档未被纳入标注的金标集的情况，即：
\begin{equation}
    \mathcal{S}_i = \hat{G}_i \setminus G_i \neq \emptyset
\end{equation}
其中 $\mathcal{S}_i$ 为第 $i$ 个实例的遗漏文档集。换言之，存在文档 $d \in \hat{G}_i$ 使得 $d \notin G_i$，即该文档客观上包含回答 $q_i$ 所需的充分信息，但未被标注者发现或标注。

成因分析。标注稀疏性是一个比标注噪声更为根本和普遍的问题，其核心成因是标注者的认知边界~\citep{zobel1998reliable, buckley2004retrieval}。在一个包含数百万甚至数千万篇文档的大规模语料库中，同一事实或知识点往往分布在多篇不同的文档中。例如，"2024年诺贝尔物理学奖得主"这一信息可能同时出现在该获奖者的个人传记页面、诺贝尔奖历史页面、相关研究领域的综述页面等多个来源中。标注者在构建评估实例时，通常只接触到其中一篇文档，并将其标注为金标。其他同样包含正确答案的文档由于未进入标注者的视野，便成为了评估中的"隐性假阴性"。

这一问题在信息检索评估的历史中早已被关注。TREC评测使用的"池化"（pooling）策略——即只对参赛系统返回的结果文档进行相关性判定——本质上也面临类似的稀疏性问题：未被任何参赛系统检索到的相关文档将永远不会被标注为相关。Buckley和Voorhees~\citep{buckley2004retrieval}的研究证明，这种不完整标注可能导致系统排名的显著变化，尤其对于那些采用了与池化阶段不同检索策略的新系统。

对评估结果的影响。标注稀疏性导致\textbf{模型性能的系统性低估}。具体而言，当检索系统成功检索到一个属于 $\mathcal{S}_i$ 的文档——即该文档确实包含了正确答案——时，由于该文档不在金标集 $G_i$ 中，评估指标将其视为一次"错误检索"（假阴性）。这导致以下后果：

（1）\textbf{Recall被系统性压低。}即使检索系统在功能意义上已经找到了用户需要的信息，但由于该文档不在 $G_i$ 中，Recall@$K$ 中的分子 $|G_i \cap \mathcal{R}_i^K|$ 不会增加，导致计算得到的召回率低于系统的实际表现。
    
（2）\textbf{排序指标被扭曲。}对于MRR和NDCG而言，如果一个属于 $\mathcal{S}_i$ 的文档被排在了排名第1位（这本应是一个优秀的检索结果），但金标文档 $d \in G_i$ 排在第5位，那么MRR将报告 $1/5$ 而非 $1/1$。NDCG也会因为在前几个位置"没有"相关文档而给予较低的分数。这不仅低估了当前检索系统的能力，更可能导致不同系统之间的排名产生误导性的逆转。
    
（3）\textbf{模型优化方向被误导。}如果使用这些带有稀疏性的标注来进行检索模型的训练（例如作为负样本），模型可能学会回避那些实际上包含正确答案但不在训练标注中的文档，从而形成对标注偏差的过拟合，而非真正提升检索能力。


标注质量的形式化度量。综合标注噪声与标注稀疏性，一个理想的高质量基准数据集应当满足：对于所有评估实例 $\mathcal{E}_i$，标注的金标集 $G_i$ 尽可能地逼近理想金标集 $\hat{G}_i$，即最小化以下两个差异：
\begin{equation}
    \min \sum_{i=1}^{M} \left( |\mathcal{N}_i| + |\mathcal{S}_i| \right) = \min \sum_{i=1}^{M} \left( |G_i \setminus \hat{G}_i| + |\hat{G}_i \setminus G_i| \right)
\end{equation}
其中 $M$ 为基准数据集中评估实例的总数。然而，精确计算 $\hat{G}_i$ 需要对整个语料库 $\mathcal{D}$ 中的每一篇文档进行相关性判定，这在大规模语料库中是不可行的。因此，本文在第四章中提出了CRB-Eval，通过设计巧妙的代理任务来间接估计 $|\mathcal{N}_i|$ 和 $|\mathcal{S}_i|$ 的严重程度，从而实现对基准质量的可操作性审计。




\section{CRB-Suite总体框架概述}
\label{sec:framework_overview}

针对上述挑战，本文提出CRB-Suite（Multi-Turn Retrieval Suite），一个集基准质量审计、高保真数据合成和检索性能评测于一体的统一框架。CRB-Suite的总体架构\footnote{具体实现细节分别在第四、五、六章中详细阐述。}包含三个紧密协作的核心模块：CRB-Eval（审计模块）、CRB-Pipeline（合成模块）和CRB-Bench（基准模块）。以下分别概述各模块的定位、核心思路及其相互关系。





\section{四维评估指标体系}

\subsection{CRB-Eval：基准质量审计模块}

模块定位。CRB-Eval是整个框架的诊断性基础设施。它的核心目标是回答一个根本性问题：\textbf{一个给定的对话检索基准数据集是否值得信赖？}——即其标注的金标集 $G_i$ 在多大程度上逼近理想金标集 $\hat{G}_i$。

核心思路。如前所述，精确计算 $\hat{G}_i$ 是不可行的。CRB-Eval的关键洞察在于：虽然无法穷举所有相关文档，但可以通过精心设计的代理任务来间接估计标注缺陷的严重程度。具体而言，CRB-Eval采用异构LLM-as-a-Judge~\citep{zheng2023judging}策略，从以下四个维度对基准数据集进行量化审计：

（1）\textbf{查询-证据对齐度}（Query-Evidence Alignment）：检测标注噪声 $\mathcal{N}_i$——判断金标文档 $d \in G_i$ 是否真正包含回答 $q_i$ 的信息。

（2）\textbf{证据完备性}（Evidence Completeness）：检测标注稀疏性 $\mathcal{S}_i$——通过可辨别性测试判断金标文档是否是语料库中最佳的答案来源。

（3）\textbf{答案-证据忠实度}（Answer-Evidence Faithfulness）：检测回答 $a_i$ 中的幻觉成分——判断回答是否完全基于金标文档集 $G_i$ 中的信息生成。

（4）\textbf{答案质量}（Answer Quality）：独立于证据文档，评估回答 $a_i$ 本身的语言质量、连贯性和有用性。


去偏机制。为确保自动化评估的可靠性，CRB-Eval整合了多重去偏设计：采用多个异构LLM的集成评分以消除单一模型的自偏好偏差；使用逐点评分（而非成对比较）以消除位置偏差；在涉及多文档的评估场景（如可辨别性测试）中，进一步对文档顺序进行随机化处理。

双重作用。CRB-Eval在框架中扮演双重角色：一方面，它作为\textbf{现有基准的审计工具}，揭示主流数据集中隐含的质量问题，为研究者理性解读评估结果提供依据；另一方面，它作为\textbf{CRB-Pipeline的质量校验器}，确保自动合成的数据达到预设的质量标准，形成闭环的质量保障机制。

\subsection{CRB-Pipeline：多智能体对话合成模块}

模块定位。CRB-Pipeline是框架的数据生产引擎。它的核心目标是：\textbf{以远低于人工标注的成本，自动合成标注质量达到甚至超越人工水平的多轮对话检索数据。}

核心思路。CRB-Pipeline将高度复杂的基准构建任务，系统性分解为三个流水线化的阶段：

（1）\textbf{知识库预处理}（Curation）：对原始语料库进行递归分块、近似去重（MinHash-LSH~\citep{broder1997resemblance}）和混合质量过滤（利用NVIDIA质量分类器和FineWeb-EDU教育价值评分器~\citep{penedo2024finewebdatasets}），确保仅高信息密度的文档段落进入后续环节。这一步从源头上保障了金标文档的质量——低质量、碎片化或重复的文档不会成为查询的答案来源。
    
（2）\textbf{贪心遍历聚类}（Clustering）：为模拟真实用户的话题浏览与切换行为，需要将语义相关的文档组织为有序序列。CRB-Pipeline提出了一种贪心遍历聚类算法，将旅行商问题（TSP）的经典贪心最近邻启发式重新应用于语义聚类场景。该算法从随机起始节点出发，迭代选择最近的未访问邻居，构建一条连续的"语义路径"，然后每 $k$ 个节点进行分割形成固定大小的文档簇。这种方法保证了：簇大小恒定（充分利用LLM上下文窗口）、每篇文档仅访问一次（消除重复偏差）、簇内语义平滑过渡（自然模拟话题流转）。
    
（3）\textbf{多智能体对话生成}（Generation）：设计三角色多智能体系统——提问者（User Simulator）负责基于文档簇和对话历史生成查询；回答者（RAG Simulator）负责严格基于指定金标文档生成回答；润色者（Refiner）负责对生成的对话进行语言自然度优化，注入指代、省略等真实对话特征。角色的解耦使得每个智能体可以专注于其特定职责，提升指令遵循度和生成质量。


成本优势。基于当前主流LLM的API定价估算，CRB-Pipeline每段对话的平均生成成本约为\TBD{单条成本美元}，相比人工标注的\$1.50至\$2.00（如Doc2Dial~\citep{feng-etal-2020-doc2dial}所报道），成本降低了约两个数量级。这种经济可行性使得大规模、跨领域的基准构建成为现实。

\subsection{CRB-Bench：对话检索评测基准模块}

模块定位。CRB-Bench是框架的输出产物和价值载体。它是利用CRB-Pipeline合成、经CRB-Eval质量验证的一个面向通用领域的大规模多轮对话检索基准。

设计原则。CRB-Bench的设计明确以真实RAG生产环境的特征为导向，在以下三个方面与现有基准形成差异化：

（1）\textbf{话题切换的真实性。}同时实现软切换（通过语义相关文档间的自然过渡）和硬切换（模拟用户在不清除对话历史的情况下突然发起全新话题）。硬切换在真实用户行为中普遍存在~\citep{spink2002multitasking}，却在现有基准中几乎缺失，是CRB-Bench区分度的重要来源之一。
    
（2）\textbf{响应特征的生产级还原。}采用与生产环境中大语言模型一致的长回复风格（平均约\TBD{平均tokens}个token，远高于多数人工标注数据集的简短回复），以及面对歧义查询时的主动推测式回答（而非学术基准中常见的澄清追问）。这使得对话历史中包含更多冗长的、需要被过滤的噪声信息，对检索模型的鲁棒性构成更严格的考验。
    
（3）\textbf{知识库的时效性与规模。}基于HuggingFace FineWeb-Edu 数据集（Common Crawl 经 LLM 教育价值过滤得到）构建约\TBD{文档规模}篇经质量筛选的文档语料库，确保语料库中包含大量预训练模型未见过的新知识，从而严格测试检索能力而非模型的参数化记忆。


与其他模块的关系。CRB-Bench并非一个静态的、一次性的产物，而是CRB-Suite框架可重复执行能力的一次实例化展示。用户可以利用CRB-Pipeline针对任意领域的私有知识库（如金融法规、医疗文献、企业内部文档等）快速生成定制化的评测基准，再利用CRB-Eval验证其质量。这种"合成-审计"的闭环能力赋予了框架强大的可扩展性和领域适应性。

\subsection{模块间的协作关系}

CRB-Suite三个模块之间的协作关系可概括为以下闭环流程：

（1）\textbf{审计驱动的问题发现。}首先使用CRB-Eval对现有基准进行审计，量化其标注噪声和标注稀疏性的严重程度。审计结果揭示了现有基准在质量和区分度方面的不足，从而验证了构建新基准的必要性。
    
（2）\textbf{合成驱动的数据生产。}基于审计发现的问题，使用CRB-Pipeline从指定知识库中自动合成新的对话检索数据。合成过程通过贪心遍历聚类和多智能体协作，从设计层面避免人工标注中认知边界导致的标注稀疏性问题。
    
（3）\textbf{审计驱动的质量保障。}合成完成后，再次使用CRB-Eval对新生成的数据进行质量审计，确保其在四个维度上均达到预设标准。未通过审计的实例将被过滤或重新生成，形成质量保障的闭环。
    
（4）\textbf{基准驱动的模型评估。}经质量验证的数据构成最终的评测基准（如CRB-Bench），用于对各类检索模型进行严格、公平的性能比较。


这种"审计$\rightarrow$合成$\rightarrow$再审计$\rightarrow$评测"的闭环设计，确保了评估结果的可靠性不再依赖于人工标注的完美性假设，而是建立在可量化、可验证的质量保障机制之上。


\section{本章小结}
\label{sec:chapter3_summary}

本章为后续的技术方案设计奠定了理论基础。首先，我们对多轮对话检索基准的核心要素进行了严格的形式化定义，包括文档语料库、对话实例、对话历史、评估实例等概念，并给出了Recall@$K$、MRR@$K$、NDCG@$K$三种主流评估指标的数学表达。随后，我们从理论层面分析了当前对话检索基准面临的两类核心缺陷——标注噪声与标注稀疏性，形式化地描述了其定义、成因及对评估结果的影响机制，指出标注噪声导致模型性能被过高估计，而标注稀疏性导致模型性能被系统性低估。最后，我们提出了CRB-Suite的总体框架设计，阐述了CRB-Eval（审计模块）、CRB-Pipeline（合成模块）和CRB-Bench（基准模块）三个核心组件的定位、核心思路及其"审计$\rightarrow$合成$\rightarrow$再审计$\rightarrow$评测"的闭环协作机制。

后续三章将分别对CRB-Eval（第四章）、CRB-Pipeline（第五章）和CRB-Bench（第六章）展开详细论述。




% \chapter{问题定义与总体框架设计}
% \label{chap:framework}

% 构建高质量的评估基准是推动检索增强生成（RAG）系统演进的先决条件。然而，在多轮对话的复杂语境下，准确量化检索性能面临着巨大的理论与工程挑战。本章将首先严格按照形式化语言定义多轮对话检索基准的核心构成，随后引入“真实证据集合”的概念，通过集合论深度剖析现有基准中存在的标注质量鸿沟，最后基于这些缺陷痛点，提出本文 CRB-Suite 的总体架构设计。

% \section{多轮对话检索基准的形式化定义}

% 在形式化层面上，一个多轮对话检索基准首先依赖于一个底层的文档语料库 $\mathcal{D} = \{d_1, d_2, \dots, d_N\}$，其中每一个 $d$ 代表一个独立的文档或段落片段。

% 基准数据集的主体由一系列评估实例构成。每一个实例对应于多轮对话中的某一个特定轮次 $i$，该轮次的交互状态可以严谨地表示为一个三元组。其简单公式表示如下：
% \begin{equation}
%     \text{Instance}_i = (H_i, q_i, G_i)
% \end{equation}

% 在该三元组中，各元素的定义与物理意义如下：
% \begin{itemize}
%     \item 当前查询 $q_i$：代表用户在第 $i$ 轮提出的当前问题，通常包含指代、省略或隐式话题依赖。
%     \item 对话历史 $H_i$：记录了当前轮次之前的所有交互上下文，表示为 $H_i = \{(q_1, a_1), \dots, (q_{i-1}, a_{i-1})\}$。$H_i$ 包含了解决 $q_i$ 所需的关键背景信息，也是对话稠密检索器（CDR）进行上下文建模的输入。
%     \item 黄金文档集合 $G_i$： 这是一个子集 $G_i \subseteq \mathcal{D}$，包含了在基准构建阶段被标注为“相关”，且足以支撑回答 $q_i$ 的标准文档集合。
% \end{itemize}

% 在评估过程中，检索模型 $f_{\text{ret}}$ 根据历史上下文 $H_i$ 和当前查询 $q_i$ 从 $\mathcal{D}$ 中召回一组文档 $D_{\text{ret}}$。传统的评估指标（如 Recall@k, NDCG@k, MRR@k）本质上是在衡量系统召回集合 $D_{\text{ret}}$ 与黄金文档集合 $G_i$ 之间的重合程度。因此，基准数据集评估结论的可靠性，绝对依赖于 $G_i$ 构建的准确性与完备性。

% \section{现有基准的核心缺陷分析}

% 尽管现有对话检索基准极大地推动了 RAG 技术的发展，但其在评估结果的可靠性与科学性上仍面临根本性挑战。评估失效的核心矛盾，在于基准数据集中的黄金文档集合 $G_i$ 与语料库 $\mathcal{D}$ 中客观存在的“真实支持文档集合” $\hat{G}_i$ 之间存在严重的质量鸿沟。

% 一个理想的高质量基准必须满足黄金文档与真实证据的完全对齐，即 $G_i \equiv \hat{G}_i$。然而，由于人类认知边界的限制或自动化启发式规则的粗糙，现有的基准数据集往往无法达到这一理想状态。基于集合的差集运算，这种质量鸿沟在评估模型时会具体表现为标注噪声和标注稀疏性等两类严重的误差，它们从不同维度扭曲了对检索系统性能的真实测量。

% \subsection{标注噪声}

% 形式化定义： 当 $G_i \setminus \hat{G}_i \neq \emptyset$ 时，即存在属于黄金文档集但不属于真实证据集的文档，此时基准数据集中便产生了标注噪声。这意味着基准的“标准答案”里混入了实际上并不相关或无法回答查询的文档片段。

% 成因解析：在多轮对话场景中，标注噪声的引入通常源于两点：其一，在自动化合成框架中，生成模型未能准确捕捉复杂对话中的话题转折，发生上下文幻觉，错误地将历史轮的支撑文档强加给当前查询（如 CORAL 中常见的问答错位）；其二，在人工标注中，标注者面对海量文本时的认知疲劳，导致其仅凭字面关键词重叠便做出轻率的相关性判断，忽略了多轮语境下细微的语义冲突。

% 评估危害：标注噪声会导致严重的模型高估现象。如果一个性能低劣的、仅基于词汇匹配的稀疏检索器恰好检索到了这些“伪相关”的噪声文档，评测指标不仅不会惩罚该模型，反而会给予其召回率上的高分。这种机制实质上奖励了错误的检索行为，使基准丧失了对具备深层语义理解能力模型的鉴别力。

% \subsection{标注稀疏性}

% 形式化定义： 当 $\hat{G}_i \setminus G_i \neq \emptyset$ 时，即客观存在于语料库中的有效支持文档未能被包含在黄金文档集合中，此时便出现了标注稀疏性。

% 成因解析： 标注稀疏性的核心成因正是“认知边界”。在面对包含数百万文档的大型语料库 $\mathcal{D}$ 时，基准构建者天然持有一种“局部视野”：他们往往是基于眼前正在阅读的特定文档逆向构造多轮对话，却无力在有限时间内全局遍历整个 $\mathcal{D}$，去寻找那些同样能够完美回答 $q_i$ 的替代文档。这些被遗漏的真实文档构成了庞大的集合 $\hat{G}_i \setminus G_i$。

% 评估危害： 标注稀疏性会导致严重的假阴性惩罚。当一个强大的神经检索器通过复杂的跨轮次推理，从知识库中挖掘出一条未被标注者发现但完全正确的证据链时，现有评估脚本会因该文档不在 $G_i$ 中而将其判定为检索失败。这种“惩罚先进”的现象意味着检索模型挖掘长尾知识的能力越强，其在稀疏基准上获得的得分反而越低，彻底摧毁了评估基准的区分度与权威性。


% \section{CRB-Suite总体框架概述}

% 为了弥合 $G_i$ 与 $\hat{G}_i$ 之间的质量鸿沟，彻底消除由于 $G_i \setminus \hat{G}_i \neq \emptyset$ 带来的过高估计，以及由于 $\hat{G}_i \setminus G_i \neq \emptyset$ 导致的假阴性惩罚，本文提出了一个集评估、合成与基准发布于一体的统一框架——CRB-Suite。

% CRB-Suite 是一套解决上述理论缺陷的系统方法论。其总体框架由三个高度协同的模块构成，形成“评估-生成-验证”的完整闭环：

% （1）CRB-Eval（基准评估模块）： 该模块直击现有基准痛点，旨在开发一套能够自动逼近 $\hat{G}_i$ 的评测工具。通过引入异构 LLM 集成裁判与四维评估指标体系（查询-证据对齐度、证据完备性、答案-证据忠实度、答案质量），CRB-Eval 能够动态检测 $G_i \setminus \hat{G}_i$（剔除噪声），并通过判别性测试主动挖掘 $\hat{G}_i \setminus G_i$（量化稀疏性），为旧基准提供科学的“去偏重打分”能力。
    
% （2）CRB-Pipeline（高保真合成模块）： 为从源头上构建高质量的 $\hat{G}_i$，本模块摒弃依赖人工局部视野的传统模式。它创新性地利用“贪心遍历聚类”在向量空间中锁定全局强相关的文档簇，确保上下文窗口的最大化利用与零冗余覆盖。随后，由提问者、回答者与润色者三智能体在受控簇内合成自然流畅的多轮交互，在确保 $G_i \approx \hat{G}_i$ 的同时，将构建成本降至人工的 $1/400$。
    
% （3）CRB-Bench（严苛验证基准）： 基于 CRB-Pipeline 输出的高质量语料，本文构建并发布 CRB-Bench。该基准显式引入硬话题切换、生产级冗长回复与长上下文歧义等真实交互现象，采用 2025 年最新知识截止与优化文档粒度。它旨在作为“压力测试台”，提供远超传统基准的模型区分度，验证现代 RAG 检索模块在复杂生产环境下的真实鲁棒性。


% 上述三大模块并非孤立运行，而是通过严密的数据流向与反馈机制构成高度协同的有机整体。在数据生成链路中，原始语料首先经递归分块、近似去重与混合质量过滤后，输入贪心遍历聚类算法构建语义连贯的文档簇；该簇作为高信息密度的上下文基底，驱动三智能体协作生成自然流畅的多轮对话，最终汇聚为 CRB-Bench。与此同时，CRB-Eval 作为贯穿始终的质量标尺，不仅对历史基准进行系统性审计，更对流水线产出的新语料执行四维量化校验，确保数据从源头到终端的严谨性。这种“以评定质、以质促生、以生验效”的闭环架构，彻底打破了传统依赖人工局部视野或静态启发式规则的构建局限，实现了数据复杂度与模型特定偏见的有效解耦，为多轮对话检索评估确立了可复用、可扩展的全局感知自动化标注新范式。


% \section{本章小结}

% 本章从数学抽象的高度出发，对多轮对话检索基准进行了严格的形式化定义 $(H_i, q_i, G_i)$。本章的核心贡献在于引入了真实证据集 $\hat{G}_i$ 的概念，并通过严谨的集合差集运算，清晰地剖析了导致现有评估失真的两大根源：导致模型高估的“标注噪声（$G_i \setminus \hat{G}_i \neq \emptyset$）”以及导致假阴性惩罚的“标注稀疏性（$\hat{G}_i \setminus G_i \neq \emptyset$）”。基于对这一质量鸿沟的深刻认知，本章系统引出了 CRB-Suite 统一框架的三大核心组件（CRB-Eval、CRB-Pipeline、CRB-Bench）及其协同逻辑。下一章将详细论述 CRB-Suite 的首个核心组件——用于量化诊断这种质量鸿沟的 CRB-Eval 评估方法及其多维指标体系。
